\subsection{Ranges}

\begin{frame}[t,fragile]{Range}
\begin{itemize}
  \item A range is any type that represents a sequence of elements.
\begin{lstlisting}
template <typename T>
concept range = requires(T& t) {
  { std::ranges::begin(t) };
  { std::ranges::end(t) };
};
\end{lstlisting}
    \begin{itemize}
      \item \cppcode{[std::begin(t), std::end(t))} is a half-open range.
      \item Both \cppcode{std::begin(t)} and \cppcode{std::end(t)} are amortized constant time.
      \item If the type of \cppcode{std::begin(t)} is a forward iterator, then \cppcode{[std::begin(t), std::end(t))} is equality preserving (i.e., multi-pass is allowed).
    \end{itemize}
\end{itemize}
\end{frame}

\begin{frame}[t,fragile]{A ranges taxonomy}
\begin{itemize}
  \item \cppcode{range}: Can be iterated from begin to end.
  \item \cppcode{output_range}: Elements can be written to the range.
  \item \cppcode{input_range}: Elements can be read from the range.
  \item \cppcode{forward_range}: Elements can be read from the range multiple times.
  \item \cppcode{bidirectional_range}: Elements can be read from the range in both directions.
  \item \cppcode{random_access_range}: Elements can be read from the range in constant time.
  \item \cppcode{contiguous_range}: All elements are stored contiguously in memory.
  \item \cppcode{sized_range}: Range with constant-time size.
\end{itemize}
\end{frame}

\begin{frame}[t,fragile]{Range based algorithms}
\begin{itemize}
  \item Ranges are the basis for a new set of algorithms in the STL.
  \item These algorithms take ranges as input and output, instead of pairs of iterators.
  \item The new algorithms are designed to be more composable and easier to use than the old ones.
\end{itemize}

\mode<presentation>{\vfill\pause}
\begin{lstlisting}
void f() {
  std::vector<double> v = get_vector();
  std::ranges::transform(v, v.begin(), [](double x) { return x * 2; });
  std::ranges::reverse(v);
  if (auto it = std::ranges::find_if(v, [](double x) { return x > 0.0; }); it != v.end()) {
    std::println("Found a positive number: {}", *it);
  }
}
\end{lstlisting}
\end{frame}